\section{Introduction}

Cyber-Physical Systems (CPS) and the so-called Internet of Things (IoT)-based access control solutions are increasingly relevant in domains such as smart homes, short-term rentals, shared facilities, and building automation. In particular, smart locks enable flexible, cloud-based management of access permissions and support automated processes such as check-in and check-out. These systems consist of multiple interacting components, including user interfaces, cloud services, databases, Application Programming Interfaces (APIs), and physical devices.

The adoption of smart locks has grown rapidly in recent years, with the global market expected to reach over US\$8 billion by 2030 \cite{refMarket}. Platforms such as Airbnb rely on smart locks to enable automated check-in and remote access management, eliminating the need for physical key exchange. As a result, access control is no longer purely physical but is increasingly mediated through software systems, APIs, and cloud services. This increases system complexity and introduces additional attack surfaces.

This increase in potential attack surfaces is accompanied by growing security concerns. Prior research has shown that Bluetooth-based smart locks can be vulnerable to eavesdropping and man-in-the-middle attacks \cite{refBluetooth}, while replay attacks can enable unauthorized unlocking by reusing captured communication \cite{diao2025masterlock}. Additionally, flaws in smart home platforms have shown that attackers can inject credentials or unlock doors via compromised cloud services \cite{refSmartHome}. These findings highlight that smart lock systems can introduce vulnerabilities across device, communication, and cloud layers.

This leads to a fundamental trade-off between usability, security, and infrastructure cost. Simple mechanisms such as static keypad codes are easy to use, but they are vulnerable to sharing, leakage, and unauthorized reuse. In contrast, stronger access control mechanisms, such as contextual validation or biometric verification, can improve protection but may also increase operational cost, Central Processing Unit (CPU) usage, memory usage, and system complexity. Similar to prior work on the measurable cost of security mechanisms, such as the cost introduced by the Hypertext Transfer Protocol Secure (HTTPS) \cite{naylor2014cost}, enhanced access control in smart-lock-based systems can also create measurable infrastructure overhead.

In this position paper, we investigate this trade-off in the context of bcube, a cloud-based booking platform for music studios. The system follows a microservice architecture and consists of a User Service, Studio Service, Booking Service, and Access Service. A Nuki Smart Lock is used as the physical access mechanism, while the cloud-based Access Service controls when a valid access code is generated or released.

We compare three access control scenarios with increasing security levels. The first scenario uses direct keypad access after a completed booking. The second scenario requires booking-code validation before a temporary keypad access code is generated. The third scenario additionally includes face verification before the access code is generated. The goal is to evaluate how these different security levels affect operational cost and infrastructure resource usage.

The evaluation focuses on three concrete metrics: cost per access attempt, average CPU utilization, and average memory utilization. These metrics are measured mainly for the Access Service, because it contains the scenario-specific access logic and is expected to show the clearest differences between the implemented security levels.

The remainder of this paper is structured as follows: Section II reviews related work, Section III presents the methodology and experimental setup, and Section IV concludes the paper.